\section{S-Matrix}


 [...]


We tipically use Heisemberg picture, so if we want to compute the amplitude of a determined process, we can write it as
\[
    \bra{f} \hat{S} \ket{i} = \bra{f, \infty} \ket{i, -\infty},
\]
where $\ket{i}$ and $\ket{f}$ are the initial and final states of the system (time independent), respectively. The S-matrix element $\bra{f} \hat{S} \ket{i}$ gives us the amplitude for the transition from the initial state $\ket{i}$ to the final state $\ket{f}$ due to the interaction described by the S-matrix.

\section{Cross Sections}

If we analyze a scattering process, such as the Rutherford scattering, we can see how a beam of particles is scattered by a target. We consider the nucleus as a still sphere, while the alpha particles are much smaller.\TODO{Add drawing} The geometric cross section of the target is given by
\[
    \sigma \equiv \pi R^2,
\]
where $R$ is the radius of the target. One can compute the number of particles that are scattered by the target, which is given by
\[
    N = n_B \cdot \sigma, \quad n_B = \frac{N_B}{A},
\]
where $n_B$ is the number density of the beam, $N_B$ is the number of particles in the beam, and $A$ is the area of the target. The cross section $\sigma$ can be interpreted as an effective area that quantifies the likelihood of a scattering event occurring when a particle from the beam interacts with the target.

Now we want to take this model and generalize it to quantum field theory, where we can compute the probability of scattering processes; we want to extend the model
\begin{itemize}
    \item for any type of scattering process, or interaction;
    \item for a probability of scattering, not just a number of particles.
\end{itemize}

If we assume a beam of particles along the $z$-axis, we can define the \textbf{impact parameter} \(\mathbf{b}\), whixch represents the perpendicular distance from the trajectory of the incoming particle to the center of the target: the coordinates perpendicular to the plane of scattering:
\[
    N = \int \mathrm{d}^2 \mathbf{b} \, n_B \, P(\mathbf{b}) = n_B \int \mathrm{d}^2 \mathbf{b} \, P(\mathbf{b}),
\]
thus we can now define the \textbf{cross section} as
\[
    \sigma = \int \mathrm{d}^2 \mathbf{b} \, P(\mathbf{b}),
\]
and its differential form as
\[
    \frac{\mathrm{d} \sigma}{\mathrm{d} \Omega} = \int \mathrm{d}^2 \mathbf{b} \, \frac{P(\mathbf{b})}{\mathrm{d} \Omega},
\]
or in the measure of the space of final states as
\[
    \frac{\mathrm{d} \sigma}{\mathrm{d} f} = \int \mathrm{d}^2 \mathbf{b} \, \frac{P(\mathbf{b})}{\mathrm{d} f}.
\]

It makes sense to use eigenstates of the momentum operator as the initial and final states of the scattering process, since they are the natural basis for describing particles in quantum field theory (errors on p ... ...)

\[
    \Delta p \Delta x \geq \frac{\hbar}{2}, \quad \Delta p = 0 \implies \Delta x \to \infty,
\]
We also know that the probability of scattering is given by the square of the amplitude \(\bra{f} \hat{S} \ket{i} \sim \delta^{(4)}(p_i - p_f)\), but then
\[
    P \sim |\bra{f} \hat{S} \ket{i}|^2 \sim |\delta^{(4)}(p_i - p_f)|^2 = \delta^{(4)}(p_i - p_f) \delta^{(4)}(0) \dots
\]

one way to overcome this problem is to consider a finite volume \(V\) and a finite time \(T\), so that the delta function can be regularized and then take only at the end the limit \(V \to \infty\) and \(T \to \infty\). This is a common technique in quantum field theory to handle divergences and make sense of scattering amplitudes. By working in a finite volume and time, we can ensure that the calculations are well-defined and then take the appropriate limits to recover physical results.

We will instead follow the Peskin and Schroeder approach, which is to consider the states as wave packets in the momentum space, which are localized in both position and momentum space, thus avoiding the issues with the delta function. This approach allows us to compute scattering amplitudes and cross sections without encountering the divergences associated with plane wave states. By using wave packets, we can ensure that the initial and final states are well-defined and that the scattering probabilities are finite and physically meaningful.

We can consider an initial state\TODO{Drawing} as a two particle state with momenta \(p_a\) and \(p_b\), and a final state as a N particle state:
[...]
For incoming states we replace \(\ket{p_i}\) with \(\ket{\phi_i}\), which is a wave packet centered around the momentum \(p_i\):
\[
    \ket{p_i} \to \ket{\phi_i} = \int {\mathrm{d}^3 \mathbf{k}}{(2 \pi)^3} \frac{1}{\sqrt{2 E_{\mathbf{k}}}} \phi_i (\mathbf{k})\ket{\mathbf{k}},
\]
with
\[
    \bra{\phi_i} \ket{\phi_i} = \int \frac{\mathrm{d}^3 \mathbf{k}}{(2 \pi)^3} |\phi_i (\mathbf{k})|^2 = 1,
\]
remembering that \(\bra{\mathbf{k}}\ket{\mathbf{k}'} = 2 E_{\mathbf{k}} (2 \pi)^3 \delta^{(3)}(\mathbf{k} - \mathbf{k}')\); all of this indicates that \(\phi_i (\mathbf{k})\) is a very narrow normalized function centered around \(p_i\).\TODO{Add graph of the wave packet}

\begin{example}[LHC collider]
    We can assume, very conservatively, that the LHC collider has \(\Delta p \sim  \SI{1}{\mega\electronvolt}\), such that ...
\end{example}

\begin{example}[Gauussian wave packet]
    We can consider a Gaussian wave packet centered around \(p_i\) with a width \(\sigma\):
    \[
        \phi (\mathbf{k}) = \dots
    \]
    This wave packet is normalized and has a well-defined momentum distribution, allowing us to compute scattering amplitudes and cross sections without encountering divergences.
        [...]
    \[
        \begin{aligned}
            \ket{p_a} \sim \ket{\phi_a} = \\
            \ket{p_b} \sim \ket{\phi_b} = \dots
        \end{aligned}
    \]
    where we exploited the dependence of \(\phi_b(\mathbf{k})\) upon \(\mathbf{b}\) the impact parameter
    \[
        \dots
    \]
\end{example}


For the outgoing states we can consider just the momentum eigenstates, since the scattering has already happened, thus we can write
\[
    \ket{\phi_i} = \ket{p_i}, \quad \forall i = 1, \dots, N.
\]


[...]

We can consider the S-matrix element for the transition by rewriting it as
\[
    \hat{S} = \mathbb{I} + i \hat{T},
\]
where the identity term corresponds to the case where no scattering occurs, while the \(i \hat{T}\) term corresponds to the case some interactions among the fields, i.e. scattering, occur. Thus we can write
\[
    \bra{p_1, \ldots,\,p_N} i \hat{T} \ket{\phi_a, \phi_b} = (2\pi)^4 \delta^{(4)}\left(p_a + p_b - \sum_{i=1}^N p_i\right) i \mathcal{M}(p_1, \ldots,\,p_N),
\]
where the delta function ensures the conservation of energy and momentum, and \(\mathcal{M}(p_1, \ldots,\,p_N)\) is the \textbf{invariant matrix element} or \textbf{scattering amplitude} for the process, which encodes the dynamics of the interaction. The invariant amplitude can be computed using Feynman diagrams and the rules of quantum field theory, and it is a key quantity in determining the cross section for the scattering process.

Thus we can now compute the differential cross section as \(\frac{\mathrm{d} \sigma}{\mathrm{d} f} = \int \mathrm{d}^2 \frac{\mathrm{d} P}{\mathrm{d} f}\), where the probability of scattering is given by
\[
    \frac{\mathrm{d} P}{\mathrm{d} f} = \vert \bra{p_1,\,\ldots,\,p_N,\,+\infty} i \hat{T} \ket{\phi_a,\, \phi_b,\, -\infty} \vert^2,
\]
and the measure of the space of final states is normalized as\footnote{We can convince ourself with the resolution of the identity for 1 particle states as ...}
\[
    \mathrm{d} f = \prod_{i=1}^N \frac{\mathrm{d}^3 \mathbf{p}_i}{(2\pi)^3} \frac{1}{2 E_{p_i}}.
\]

Now we can put everything back together and write the final expression for the differential cross section after ......... conti

\[
    \Phi = 2 E_a 2 E_b \vert v_a - v_b \vert = 4 \vert E_b p_a^z - E_a p_b^z \vert.
\]

\(\Phi\) is invariant under boosts along the beam axis, and rotations around the beam axis, thus it can only depend on the Lorentz invariants of the process, and thus it transforms as the inverse of an area in the plane perpendicular to the beam axis, thus \dots

We can also perform a change of reference frame in order to reverse the problem: we can consider the target as the beam and the beam as the target, it all depends on the symmetry of the problem, and the choice of the reference frame. This is a common technique in scattering theory, where we can choose a reference frame that simplifies the calculations (even center of mass, it depends) and makes it easier to analyze the scattering process. By considering different reference frames, we can gain insights into the dynamics of the interaction and compute cross sections more efficiently.

\begin{example}[2 \(\to\) 2 scattering in center of mass frame]\TODO{Add drawing}
    We can consider a simple scattering process where two particles collide and produce two particles in the final state. In the center of mass frame, the initial momenta of the particles are equal and opposite, and we can analyze the scattering process using the invariant amplitude \(\mathcal{M}\) and the phase space factor \(\Phi\) to compute the differential cross section for this process.

    The incoming particles have momenta \(p_1\) and \(p_2\), and the outgoing particles have momenta \(p_3\) and \(p_4\), and their spatial parts is conserved: \(\mathbf{p}_1 + \mathbf{p}_2 = \mathbf{p}_3 + \mathbf{p}_4 = 0\) (since we are in the center of mass frame). The energy of the particles is given by the energy momentum relation \(E_i = \sqrt{\mathbf{p}_i^2 + m_i^2}\), where \(m_i\) is the mass of the particle. We can write
    \[
        \mathrm{d} \Pi_2 = \frac{\mathrm{d}^3 \mathbf{p}_3}{(2\pi)^3 2 E_3} \frac{\mathrm{d}^3 \mathbf{p}_4}{(2\pi)^3 2 E_4} \delta^{(4)}(p_1 + p_2 - p_3 - p_4),
    \]
    which, after integrating in \(\mathrm{d}^3 \mathbf{p}_4\) using the delta function, gives \(\mathbf{p}_4 = - \mathbf{p}_3\), and thus using then polar coordinates for \(\mathrm{d}^3 \mathbf{p}_3 = p^2 \mathrm{d} p \mathrm{d} \Omega\), we can write
    \[
        E_3 = \sqrt{p^2 + m_3^2}, \quad E_4 = \sqrt{p^2 + m_4^2},
    \]
    and thus
    \[
        \int \mathrm{d} p \delta(E_1 + E_2 - E_3 - E_4) = \frac{1}{\left\vert \frac{p}{E_3} + \frac{p}{E_4} \right\vert} = \frac{E_3 E_4}{p (E_3 + E_4)} = \frac{E_3 E_4}{p E},
    \]
    where the last energy \(E\) is the total energy of the system in the center of mass frame, which is given by \(E = E_1 + E_2 = E_3 + E_4\). Thus we can write
    \[
        \mathrm{d} \Pi_2 = \mathrm{d} \Omega \frac{p}{16 \pi^2 E},
    \]
    with a flux \(\Phi = 2 E_a 2 E_b \vert v_a - v_b \vert\), with now
    \[
        \vert v_a - v_b \vert = \left\vert \frac{\vert \mathbf{p}_1 \vert }{E_1} + \frac{\vert \mathbf{p}_2 \vert }{E_2} \right\vert = p_{in} \frac{E}{E_1 \, E_2},
    \]
    where \(p_{in}\) is the magnitude of the incoming momenta (which are equal in the center of mass frame) \(p_{in} = \vert \mathbf{p}_1 \vert = \vert \mathbf{p}_2 \vert \). Thus we can write the final expression for the differential cross section as
    \[
        \frac{\mathrm{d} \sigma}{\mathrm{d} \Omega} = \frac{1}{64 \pi^2 E^2} \frac{p}{p_{in}} |\mathcal{M}|^2,
    \]
    where we used the fact that \(E_1 E_2 = \frac{E^2}{4}\) in the center of mass frame. This expression allows us to compute the differential cross section for the 2 \(\to\) 2 scattering process in the center of mass frame, given the invariant amplitude \(\mathcal{M}\) and the kinematic variables of the process.
    \begin{remark}
        Most of the time this problem is independent of the azimuthal angle \(\phi\), thus we can integrate over it and write \(\mathrm{d}\Omega = 2\pi \mathrm{d} \cos \theta\), thus simplifying the expression for the differential cross section.
    \end{remark}
\end{example}


\section{Decay Rates}

We can similarly deal with \textbf{unstable particles}, which are particles that can decay into other particles. The decay rate of an unstable particle is given by the probability per unit time that the particle will decay, and it can be computed using the S-matrix formalism in a similar way to how we compute scattering amplitudes and cross sections for stable particles.\TODO{Add drawing.}
By analyzing the transition amplitudes for the decay process, we can determine the decay rate and the branching ratios for different decay channels, which provide insights into the properties of the unstable particle and its interactions with other particles.

The decay rate \(\Gamma\) of an unstable particle can be computed as
\[
    \Gamma = \frac{\text{Number of decays per unit time}}{\text{Number of particles}},
\]
and we can define the lifetime \(\tau\) of the particle as the inverse of the decay rate:
\[
    \tau = \frac{1}{\Gamma}.
\]
The decay rate can be computed using the S-matrix element for the transition from the initial state (the unstable particle) to the final state (the decay products), and it is related to the invariant amplitude \(\mathcal{M}\) for the decay process. We cannot directly apply the limit for \(t_i \to -\infty\), since the particle is unstable and thus it cannot exist for an infinite amount of time, thus we need to consider a finite time interval for the decay process, but we can make a simplification: if we assume the particle to have a lifetime way longer than the time scale of the interaction, we can still apply the S-matrix formalism and compute the decay rate using the invariant amplitude for the decay process, while treating the unstable particle as if it were stable during the interaction.

Defining the differential decay rate as
\[
    \frac{\mathrm{d} \Gamma}{\mathrm{d} f} = \frac{1}{T} \mathrm{d} P,
\]
where \(T\) is the time interval during which the decay process occurs \(T = t_f-t_i\), and \(\mathrm{d} P\) is the probability of the decay process occurring within that time interval, we can compute the total decay rate by integrating over the phase space of the final states; let us first write the expression for t delta function in the limit of finite and large \(T\):
\[
    \delta(0) = \int_{t_i}^{t_f} \mathrm{d} t \, e^{i (E_i - E_f) t} |_{t_f \to \infty, t_i \to -\infty},
\]

After surpassing the subtleties of defining the asymptotic states for unstable particles, we can write the expression for the decay rate after integrating over the phase space of the final states as
\[
    \mathrm{d} \Gamma = \frac{1}{2 m} |\mathcal{M}|^2 \mathrm{d} \Pi_N,
\]
where \(m\) is the mass of the unstable particle, \(\mathcal{M}\) is the invariant amplitude for the decay process, and \(\mathrm{d} \Pi_N\) is the phase space factor for the N-particle final state, which accounts for the kinematics of the decay products and ensures that energy and momentum are conserved in the decay process. This expression allows us to compute the decay rate for an unstable particle given the invariant amplitude and the kinematic variables of the decay process.

\paragraph{Identical particles.} When dealing with identical particles among the decay products, we need to take into account the symmetrization (for bosons) or antisymmetrization (for fermions) of the wave function. This means that the scattering amplitude must be modified to include the appropriate symmetry factors, which can affect the calculation of cross sections and decay rates in order to consider the spin statistics theorem
\[
    \ket{p_3,\,p_4} = \pm \ket{p_4,\,p_3}.
\]
This can be accomplished in two ways: one can restrict the integration over the phase space to a specific region that accounts for the indistinguishability of the particles (thus over all final inequivalent states), or one can include a symmetry factor in the calculation of the amplitude (such as the permutation number \(\frac{1}{N!}\)) to account for the number of identical particles in the final state. Both approaches ensure that the physical predictions are consistent with the principles of quantum mechanics and the statistics of identical particles.